\chapter{دوال \(L\) الآلية، دون التحدب، ومدخل إلى لانجلاندز}
\label{ch:automorphic-l-subconvexity-langlands}

\section*{نطاق الفصل وحالته}

ينقل هذا الفصل البنية التي ظهرت في الفصل العشرين من معاملات هيكه
والأشكال الهولومورفية وأشكال مااس إلى دوال \(L\) الآلية. النواة
التعليمية هي \(GL(2)/\mathbb Q\): نثبت قاموس التطبيع وحاصل أويلر في
الحالة غير المتشعبة، ثم نعرض الاستمرار التحليلي والمعادلة الوظيفية
والموصل التحليلي وحد التحدب، ونشرح معنى كسر حاجز التحدب. نختم بعرض
محدود للمنظور المحلي--العالمي وبرنامج لانجلاندز.

\noindent\textbf{حالة الفصل:}
\textenglish{AUTHORED-DRAFT / NON-CITABLE}. اجتاز نطاق الفصل وسجل
أدلته وجدول تطبيعاته وخريطة برهانه المراجعة المستقلة قبل التأليف. لا
تصبح نتائجه قابلة للاستشهاد قبل البناء والتدقيق والمراجعة اللاحقة
واعتماد مالك المشروع.

\subsection*{حراس النطاق}

\begin{itemize}
  \item خط التناظر هو \(\Re s=1/2\)، ومعاملات هيكه مطبعة تحليليًا.
  \item الصيغ التفصيلية لعوامل غاما تخص الأشكال الجديدة على
  \(GL(2)/\mathbb Q\).
  \item مبرهنة Michel--Venkatesh تذكر فوق حقل أعداد ثابت، ولا يثبت
  الفصل برهانها.
  \item برنامج لانجلاندز يعرض إطارًا تاريخيًا ومفاهيميًا؛ لا ندعي أن
  \textenglish{functoriality} العامة مبرهنة.
  \item دوال Rankin--Selberg و\(GL(n)\) العام خارج النطاق.
\end{itemize}

\subsection*{نواتج التعلم}

بنهاية الفصل يستطيع القارئ أن:
\begin{enumerate}
  \item يحول معاملات فورييه الخام لشكل هولومورفي إلى معاملات هيكه
  ذات خط تناظر \(1/2\).
  \item يشتق العامل الأويلري من علاقة هيكه عند القوى الأولية.
  \item يميز العامل المحلي غير المتشعب عن العوامل المتشعبة
  والأرخميدية.
  \item يكتب الدالة المكتملة والمعادلة الوظيفية مع التمثيل المقابل.
  \item يحسب الموصل التحليلي تقريبيًا في حالتي الهولومورفي ومااس.
  \item يفرق بين حد التحدب وحد ليندلوف وحد دون التحدب.
  \item يحدد بدقة ما توحّده مبرهنة Michel--Venkatesh وما تثبته.
  \item يقرأ حاصل أويلر بوصفه تجميعًا عالميًا لبيانات محلية.
  \item يفصل النتائج المثبتة عن الإطار الحدسي في مدخل لانجلاندز.
\end{enumerate}

\section{من معاملات هيكه إلى السلسلة المعيارية}

ليكن
\[
f(z)=\sum_{n\ge1}a_f(n)e(nz)
\]
شكلًا جديدًا هولومورفيًا أوليًا من الوزن \(k\)، مطبعًا بـ
\(a_f(1)=1\). السلسلة الكلاسيكية
\[
\sum_{n\ge1}a_f(n)n^{-s}
\]
لها خط تناظر عند \(\Re s=k/2\). لكن لغة دوال \(L\) الآلية توحّد
خطوط التناظر بنقلها إلى \(\Re s=1/2\).

\begin{definition}[التطبيع التحليلي لمعاملات هيكه]
\label{def:ch21-hecke-normalization}
\resultid{ANT-DEF-21-01}
\provedhere
نضع
\[
\lambda_f(n)=a_f(n)n^{-(k-1)/2},
\qquad
L(s,f)=\sum_{n\ge1}\frac{\lambda_f(n)}{n^s}.
\]
وعندئذ
\[
L(s,f)=L_{\mathrm{cl}}\!\left(s+\frac{k-1}{2},f\right),
\]
حيث \(L_{\mathrm{cl}}(w,f)=\sum a_f(n)n^{-w}\).
\end{definition}

التحويل جبري، لكنه ليس تجميليًا: المعادلة الكلاسيكية التي تربط \(w\)
بـ\(k-w\) تصبح بعد وضع \(w=s+(k-1)/2\) معادلة تربط \(s\) بـ
\(1-s\). وهكذا يصبح الخط الحرج نفسه في دوال ديريشليه والأشكال
الهولومورفية وأشكال مااس.

\begin{remark}[مااس: تطبيع الشكل وتطبيع دالة \(L\)]
في الفصل العشرين ميزنا بين \(\|u\|_2=1\) وبين تطبيع معامل هيكه الأول.
دالة \(L(s,u)\) القياسية تستعمل معاملات هيكه \(\lambda_u(n)\) مطبعة
بـ\(\lambda_u(1)=1\). تغيير حجم الدالة الذاتية \(u\) لا ينبغي أن
يغير دالة \(L\) القياسية بعد هذا التطبيع.
\end{remark}

\section{حاصل أويلر من علاقات هيكه}

في نواة المستوى \(1\) تحقق معاملات الشكل الذاتي العلاقة المسجلة في
الفصل العشرين:
\[
\lambda_f(m)\lambda_f(n)
=
\sum_{d\mid(m,n)}
\lambda_f\!\left(\frac{mn}{d^2}\right).
\]
وبخاصة، عند قوة أولية:
\[
\lambda_f(p^{r+1})
=
\lambda_f(p)\lambda_f(p^r)-\lambda_f(p^{r-1}),
\qquad r\ge1.
\]

\begin{proposition}[حاصل أويلر في النواة غير المتشعبة]
\label{prop:ch21-euler-product}
\resultid{ANT-PROP-21-01}
\provedhere
في نصف مستوى التقارب المطلق، ولشكل المستوى \(1\) السابق، لدينا
\[
L(s,f)
=
\prod_p
\left(1-\lambda_f(p)p^{-s}+p^{-2s}\right)^{-1}.
\]
وإذا عرّفنا عددي ساتاكي \(\alpha_f(p),\beta_f(p)\) بجذري
\[
X^2-\lambda_f(p)X+1=0,
\]
فإن
\[
L_p(s,f)
=
\frac{1}{(1-\alpha_f(p)p^{-s})(1-\beta_f(p)p^{-s})}.
\]
\end{proposition}

\begin{proof}
ضع
\[
G_p(X)=\sum_{r\ge0}\lambda_f(p^r)X^r.
\]
نضرب العلاقة التراجعية في \(X^{r+1}\) ونجمع على \(r\ge1\). مع
\(\lambda_f(1)=1\) نحصل على
\[
(1-\lambda_f(p)X+X^2)G_p(X)=1.
\]
إذن
\[
G_p(X)=\frac1{1-\lambda_f(p)X+X^2}.
\]
ومن الضربية عند الأعداد المتباينة أوليًا، ومن التقارب المطلق الذي يجيز
إعادة الترتيب، ينتج حاصل أويلر بعد وضع \(X=p^{-s}\). والتحليل إلى
عاملين خطيين هو تعريف \(\alpha_f(p),\beta_f(p)\).
\end{proof}

\begin{remark}[لا نفترض رامانوجان سرًا]
العلاقة \(\alpha_f(p)\beta_f(p)=1\) تأتي من كثير الحدود في حالة المستوى
\(1\)، لكنها لا تعطي وحدها
\(|\alpha_f(p)|=|\beta_f(p)|=1\). هذه الأخيرة نتيجة عميقة في الحالة
الهولومورفية، وحدسية عامة في صيغ أوسع. بناء حاصل أويلر لا يحتاج إلى
إدخالها فرضًا غير معلن.
\end{remark}

\section{العوامل المحلية}

الجداء السابق يقترح أن دالة \(L\) العالمية تتكون من عامل عند كل موضع:
الأوليات المنتهية والموضع الأرخميدي. في اللغة التمثيلية نكتب تمثيلًا
آليًا عالميًا على الصورة المقيدة
\[
\pi\simeq\bigotimes_v'\pi_v.
\]
عند معظم الأوليات يكون \(\pi_p\) غير متشعب، فتكون بياناته ممثلة بزوج
ساتاكي.

\begin{definition}[العامل المحلي غير المتشعب]
\label{def:ch21-unramified-local-factor}
\resultid{ANT-DEF-21-02}
\citedresult[\cite{JacquetLanglands1970}، Ch.~III، \S11]
إذا كان \(\pi_p\) تمثيلًا غير متشعب لـ\(GL(2,\mathbb Q_p)\) ذي
معلمتي ساتاكي \(\alpha_{\pi,1}(p),\alpha_{\pi,2}(p)\)، فنضع
\[
L_p(s,\pi)
=
\prod_{j=1}^{2}
\left(1-\alpha_{\pi,j}(p)p^{-s}\right)^{-1}.
\]
والجداء العالمي المنتهي هو
\[
L(s,\pi)=\prod_{p<\infty}L_p(s,\pi)
\]
حيث يتقارب مطلقًا أولًا في نصف مستوى أيمن.
\end{definition}

عند الأوليات التي تقسم المستوى تتغير درجة العامل أو معلماته بحسب النوع
المحلي؛ لذلك لا يجوز نسخ العامل غير المتشعب إليها آليًا. أما العامل
الأرخميدي فيسجل الوزن أو المعلمة الطيفية.

نستعمل الاصطلاحين
\[
\Gamma_{\mathbb R}(s)=\pi^{-s/2}\Gamma(s/2),
\qquad
\Gamma_{\mathbb C}(s)=2(2\pi)^{-s}\Gamma(s),
\]
مع الهوية
\[
\Gamma_{\mathbb C}(s)
=
\Gamma_{\mathbb R}(s)\Gamma_{\mathbb R}(s+1).
\]
في الحالتين التعليميتين:
\[
L_\infty(s,f)
=
\Gamma_{\mathbb C}\!\left(s+\frac{k-1}{2}\right)
\]
للشكل الهولومورفي، و
\[
L_\infty(s,u)
=
\Gamma_{\mathbb R}(s+\kappa+ir)
\Gamma_{\mathbb R}(s+\kappa-ir)
\]
لشكل مااس من الوزن صفر، ذي قيمة لابلاس \(1/4+r^2\) وزوجية
\(\kappa\in\{0,1\}\).

\section{الدالة المكتملة والمعادلة الوظيفية}

ليكن \(N_\pi\) الموصل الحسابي. نضع، ضمن التطبيع المجمد،
\[
\Lambda(s,\pi)
=
N_\pi^{s/2}L_\infty(s,\pi)L(s,\pi).
\]
قد تختلف بعض المراجع بعامل لا يعتمد على \(s\)، ولا يغير ذلك خط التناظر
أو مضمون المعادلة الوظيفية إذا استعمل الاصطلاح نفسه في الطرفين.

\begin{theorem}[الاستمرار التحليلي والمعادلة الوظيفية لـ\(GL(2)\)]
\label{thm:ch21-functional-equation}
\resultid{ANT-THM-21-01}
\citedresult[\cite{JacquetLanglands1970}، Ch.~III، \S11، Thm.~11.1]
لتمثيل آلي حدبي \(\pi\) على \(GL(2,\mathbb A_F)\)، تمتد دالة
\(L(s,\pi)\) المكتملة إلى دالة تامة، وتضبط في الشرائط الرأسية، وتحقق
\[
\Lambda(s,\pi)
=
\varepsilon(\pi)
\Lambda(1-s,\widetilde\pi),
\qquad
|\varepsilon(\pi)|=1
\]
بعد اختيار التطبيع الوحدوي المتوافق، حيث \(\widetilde\pi\) التمثيل
المقابل.
\end{theorem}

\begin{remark}[لماذا يظهر التمثيل المقابل؟]
المعاملات المحلية للتمثيل المقابل تعكس معلمات ساتاكي، والمعادلة
الوظيفية تنقل البيانات من \(s\) إلى \(1-s\). لا يجوز حذف
\(\widetilde\pi\) إلا إذا كانت \(\pi\) ذاتية الازدواج، أو بعد بيان
التطابق المناسب.
\end{remark}

\begin{remark}[الحالة غير الحدبية]
التمثيلات المتولدة من سلاسل Eisenstein قد تعطي دوال ميرومورفية وأقطابًا.
صياغة المبرهنة السابقة قصدت الحالة الحدبية، فلا نعمم تمامها إلى كل طيف
آلي دون قيد.
\end{remark}

\section{الموصل الحسابي والموصل التحليلي}

الموصل الحسابي \(N_\pi\) يقيس التشعب عند المواضع المنتهية، لكنه لا
يرى كبر الوزن أو المعلمة الطيفية أو ارتفاع النقطة \(t\). لذلك نحتاج
معيارًا يجمع التعقيد المنتهي والأرخميدي.

\begin{definition}[الموصل التحليلي]
\label{def:ch21-analytic-conductor}
\resultid{ANT-DEF-21-03}
\citedresult[\cite{MichelVenkatesh2010}، \S3.1.8؛ \cite{IwaniecSarnak2000}]
إذا كتب العامل الأرخميدي عند \(\mathbb Q\) على الصورة
\[
L_\infty(s,\pi)=
\prod_{j=1}^{2}\Gamma_{\mathbb R}(s+\mu_j),
\]
فنستعمل
\[
C(\pi,t)
=
N_\pi\prod_{j=1}^{2}(2+|\mu_j+it|),
\qquad
C(\pi)=C(\pi,0).
\]
وبصورة تمثيلية مكافئة نكتب
\[
C(\pi,t)=C\!\left(\pi\otimes|\det|^{it}\right).
\]
التعريفات القياسية التي تستبدل \(2+|z|\) بتعبير مقارن تعد متكافئة
حتى ثوابت تعتمد على الدرجة والتطبيع.
\end{definition}

في الحالة الهولومورفية نحصل على المقارنة
\[
C(\pi,t)\asymp N_\pi(k+|t|)^2.
\]
وفي حالة مااس:
\[
C(\pi,t)
\asymp
N_\pi(2+|t+r|)(2+|t-r|).
\]
الصيغة الثانية تكشف ظاهرة هبوط الموصل عندما تتقارب معلمات معينة؛ ولهذا
قد يكون حد قوي في \(|t|+|r|\) غير دون محدب قرب مناطق الهبوط.

\section{حد التحدب}

المعادلة الوظيفية تعطي معلومات على جانبي الشريط الحرج. ومع تقديرات
النمو ومبدأ فراجمن--ليندلوف يمكن نقل تقديرات الحافتين إلى الداخل. لكن
الفصل الثالث من الموسوعة لم يثبت النسخة اللازمة من هذا المبدأ؛ لذلك
نقتبس النتيجة القياسية ولا نخفيها في برهان صوري.

\begin{theorem}[حد التحدب]
\label{thm:ch21-convexity-bound}
\resultid{ANT-THM-21-02}
\citedresult[\cite{MichelVenkatesh2010}، \S1.1؛ \cite{IwaniecSarnak2000}]
لأي \(\varepsilon>0\)، وفي عائلة من دوال \(L\) ذات درجة وحقل
مثبتين وتطبيع موحد، لدينا عند المركز
\[
L(1/2,\pi)
\ll_{n,F,\varepsilon}
C(\pi)^{1/4+\varepsilon}.
\]
وبالالتواء الأرخميدي:
\[
L(1/2+it,\pi)
\ll_{n,F,\varepsilon}
C(\pi,t)^{1/4+\varepsilon}.
\]
\end{theorem}

الأس \(1/4\) ليس تخمينًا للحجم الحقيقي؛ إنه حاجز عام ناتج عن
الاستيفاء بين منطقتي التحكم. تخمين ليندلوف يتنبأ، في اللغة نفسها، بأن
\[
L(1/2,\pi)\ll_{\varepsilon}C(\pi)^\varepsilon.
\]
إذن يوجد فضاء واسع بين ما تمنحه الأدوات العامة وما يتوقع أن يكون
صحيحًا.

\section{ما معنى دون التحدب؟}

لا تكفي عبارة «أفضل من حد تافه». الاصطلاح يقارن مباشرة بالأس
\(1/4\) في الموصل التحليلي الكلي.

\begin{definition}[مسألة دون التحدب]
\label{def:ch21-subconvexity}
\resultid{ANT-DEF-21-04}
\provedhere
لتكن \(\mathcal F\) عائلة محددة من التمثيلات الآلية. نقول إن حدًا
دون محدب بقوة ثابتة مثبت على \(\mathcal F\) إذا وجد \(\delta>0\)
بحيث
\[
L(1/2,\pi)
\ll_{\mathcal F}
C(\pi)^{1/4-\delta}
\]
لكل \(\pi\in\mathcal F\)، مع التصريح بما يسمح له أن يتغير وباعتماد
الثابت الضمني. وينطبق التعريف على جانب \(t\) باستبدال \(\pi\) بـ
\(\pi\otimes|\det|^{it}\).
\end{definition}

\begin{remark}[ثلاثة جوانب لا تختلط]
\begin{enumerate}
  \item \textenglish{\(t\)-aspect}: الشكل ثابت و\(|t|\to\infty\).
  \item جانب المستوى أو الموصل: تتغير البيانات المنتهية.
  \item الجانب الطيفي: تتغير معلمة لابلاس أو الوزن.
\end{enumerate}
قد يكون التقدير دون محدب في جانب، وغير موحد أو حتى غير دون محدب في
منطقة تجمع جانبين. الموصل التحليلي هو اللغة التي تكشف ذلك.
\end{remark}

\begin{remark}[دون التحدب الضعيف]
ادخار قوة من \(C\) أقوى من ادخار لوغاريتمي. لا نساوي بين
\(C^{1/4}/(\log C)^A\) وبين \(C^{1/4-\delta}\). هذا الفصل يستعمل
مصطلح دون التحدب في معناه ذي الادخار بقوة ما لم يصرح بخلافه.
\end{remark}

\section{مبرهنة Michel--Venkatesh}

الإنجاز المركزي في النطاق الحالي أنه لا يقتصر على جانب واحد منفصل، بل
يصاغ بدلالة الموصل التحليلي عبر الجوانب المختلفة.

\begin{theorem}[Michel--Venkatesh]
\label{thm:ch21-michel-venkatesh}
\resultid{ANT-THM-21-03}
\citedresult[\cite{MichelVenkatesh2010}، Thm.~1.1]
ليكن \(F\) حقل أعداد ثابتًا، و\(\mathbb A_F\) حلقة أديلاته. يوجد
ثابت مطلق \(\delta>0\) بحيث إذا كانت \(\pi\) تمثيلًا آليًا على
\(GL_1(\mathbb A_F)\) أو \(GL_2(\mathbb A_F)\)، ذا طابع مركزي
وحدوي، فإن
\[
L(1/2,\pi)
\ll_F
C(\pi)^{1/4-\delta}.
\]
\end{theorem}

\begin{remark}[حدود الصياغة]
\begin{itemize}
  \item \(\delta\) مطلق، لكن الثابت الضمني يعتمد على الحقل الثابت
  \(F\).
  \item لا نثبت المبرهنة هنا، ولا ننسب إلى الورقة قيمة عددية محسنة
  لـ\(\delta\).
  \item تدخل القيم \(L(1/2+it,\pi)\) في المبرهنة عبر الالتواء
  بـ\(|\det|^{it}\).
  \item لا نخلط هذه الصيغة بـTheorem 1.2 في الورقة، المتعلقة بدوال
  Rankin--Selberg لتمثيلين على \(GL(2)\).
\end{itemize}
\end{remark}

تأتي قوة النتيجة من صياغة تحليلية وتمثيلية موحدة. البرهان يستعمل علاقات
فترات آلية مطبعة قانونيًا، والتضخيم، والتحليل المحلي، ولا يختزل إلى
تطبيق منفرد لصيغة Kuznetsov أو لمجاميع أسية. لهذا أبقيناها نتيجة
مقتبسة كاملة بدل تقديم «برهان تخطيطي» قد يخفي المدخلات العميقة.

\section{التمثيل العالمي وبياناته المحلية}

\begin{definition}[التمثيل الآلي العالمي وعوامله المحلية]
\label{def:ch21-global-automorphic-representation}
\resultid{ANT-DEF-21-05}
\citedresult[\cite{JacquetLanglands1970}؛ \cite{LanglandsEulerProducts1971}]
في نطاق هذا الفصل، التمثيل الآلي العالمي \(\pi\) لـ
\(GL(2,\mathbb A_F)\) هو مكوّن غير قابل للاختزال من الطيف الآلي،
وله تحليل جداء مقيد
\[
\pi\simeq\bigotimes_v'\pi_v.
\]
ترتبط بكل \(\pi_v\) بيانات محلية وعامل \(L_v(s,\pi_v)\). ووفق
اصطلاح هذا الفصل نجمع المواضع المنتهية في
\[
L(s,\pi)=\prod_{v\nmid\infty}L_v(s,\pi_v),
\]
ثم نضم العوامل الأرخميدية والموصل الحسابي في الدالة المكتملة
\(\Lambda(s,\pi)\).
\end{definition}

هذه اللغة تشرح لماذا لا يكون حاصل أويلر مجرد حيلة لتوليد سلسلة
ديريشليه: كل أولي يرى تمثيلًا محليًا، والدالة العالمية تحفظ توافق هذه
المعلومات. عند معظم المواضع تكون البيانات غير متشعبة وبسيطة؛ الصعوبة
تتركز في المواضع المتشعبة والأرخميدية وفي شرط أن تأتي العوامل جميعًا
من كائن عالمي واحد.

\section{مدخل منضبط إلى برنامج لانجلاندز}

بدأ المنظور الذي طوره لانجلاندز من سؤال توحيدي: لماذا تتشابه دوال
\(L\) الآتية من الشخصيات، والأشكال المعيارية، والتمثيلات الغالوية؟
الجواب المقترح ليس مساواة دوال منعزلة، بل شبكة مراسلات بين تمثيلات
مجموعات مختلفة تحفظ العوامل المحلية ودوال \(L\) والمعادلات الوظيفية
\cite{LanglandsProblems1970,LanglandsEulerProducts1971}.

أبسط جسر في موسوعتنا هو:
\[
\text{شخصيات ديريشليه/هيكه}
\longleftrightarrow
GL(1),
\]
ثم
\[
\text{الأشكال الجديدة الهولومورفية أو مااس}
\longleftrightarrow
\text{تمثيلات آلية على }GL(2).
\]
هذا لا يثبت البرنامج العام، لكنه يوضح لماذا بدت موضوعات الفصلين 7 و20
حالتين من لغة واحدة.

\subsection*{مبدأ النقل الدالي}

إذا كان لدينا تمثيل آلي لمجموعة ما وتمثيل مناسب لمجموعة لانجلاندز
المقابلة، يتنبأ مبدأ \textenglish{functoriality} بنقل آلي إلى مجموعة
أخرى بحيث تتوافق المعلمات المحلية، ومن ثم تتوافق دوال \(L\) المناسبة.
في حالات خاصة عميقة ثبتت عمليات رفع ونقل، لكن الصورة العامة مفتوحة.

\begin{remark}[مسألة مفتوحة: النقل الدالي العام]
\label{open:ch21-general-functoriality}
\resultid{ANT-OPEN-21-01}
\deferredresult{برنامج لانجلاندز العام}
إثبات \textenglish{functoriality} بالمستوى العام الذي يقترحه برنامج
لانجلاندز مسألة مفتوحة. لا تسجل هذه الفقرة مبرهنة، ولا تدعي أن توافقًا
شكليًا بين عوامل أويلر ينشئ تلقائيًا تمثيلًا آليًا عالميًا.
\end{remark}

\subsection*{ثلاث طبقات معرفية}

ينبغي للقارئ فصل:
\begin{enumerate}
  \item \textbf{نتائج مثبتة:} مثل نظرية هيكه لـ\(GL(2)\) ومبرهنة
  Michel--Venkatesh في نطاقها.
  \item \textbf{مراسلات أو عمليات نقل مثبتة في حالات خاصة:} لا تعمم
  خارج فرضياتها.
  \item \textbf{إطار حدسي عام:} يوجه البحث، لكنه ليس مخزنًا لمبرهنات
  جاهزة.
\end{enumerate}
هذا الفصل يكتفي بهذه الخريطة، ولا يدخل في تعريف \(L\)-groups أو
المراسلة المحلية العامة.

\section{كيف تتصل الفصول السابقة بهذا الفصل؟}

\begin{center}
\begin{tabular}{lll}
\textbf{المصدر} & \textbf{المعلومة الواردة} & \textbf{دورها هنا}\\
الفصل 7 & دوال ديريشليه والشخصيات & نموذج \(GL(1)\)\\
الفصل 18 & المجاميع الأسية & أدوات في مسائل خاصة، لا برهان عام\\
الفصل 20 & هيكه ومااس وصيغ التتبع & مدخلات \(GL(2)\) والطيف\\
الفصل 21 & الموصل ودون التحدب & توحيد الحجم عبر العائلات
\end{tabular}
\end{center}

صيغة Petersson أو Kuznetsov قد تحول متوسطات معاملات إلى مجاميع
Kloosterman، وهذا مفيد في براهين دون التحدب لبعض العائلات. لكن مبرهنة
Michel--Venkatesh العامة لا تتبع منطقيًا من مجرد وجود صيغة تتبع في
الفصل العشرين. الفصل بين «أداة ممكنة» و«اعتماد برهاني» يمنع الدور.

\section{أخطاء شائعة}

\begin{enumerate}
  \item استعمال \(a_f(n)\) الخام داخل دالة مركزها \(1/2\).
  \item نقل عامل غاما الهولومورفي إلى شكل مااس.
  \item كتابة \(N_\pi=C(\pi,t)\) وإهمال المعلمات الأرخميدية.
  \item حذف \(\widetilde\pi\) من المعادلة الوظيفية بلا ذاتية ازدواج.
  \item وصف أي ادخار لوغاريتمي بأنه ادخار بقوة.
  \item قول «موحد في جميع الجوانب» دون تثبيت الحقل وبيان الثابت.
  \item افتراض رامانوجان عند تعريف معلمات ساتاكي.
  \item الخلط بين Theorem 1.1 وTheorem 1.2 في Michel--Venkatesh.
  \item عرض برنامج لانجلاندز العام بوصفه مبرهنة مكتملة.
\end{enumerate}

\section{تمارين}

\subsection{المستوى الأول}
\begin{enumerate}
  \item تحقق من أن التحويل \(w=s+(k-1)/2\) ينقل \(w\leftrightarrow k-w\)
  إلى \(s\leftrightarrow1-s\).
  \item استخرج الحدود الأربعة الأولى من
  \((1-\lambda_f(p)X+X^2)^{-1}\).
  \item أثبت الهوية
  \(\Gamma_{\mathbb C}(s)=\Gamma_{\mathbb R}(s)
  \Gamma_{\mathbb R}(s+1)\) من صيغة مضاعفة غاما.
  \item احسب رتبة \(C(\pi,t)\) لشكل هولومورفي ثابت الوزن عندما
  \(|t|\to\infty\).
\end{enumerate}

\subsection{المستوى الثاني}
\begin{enumerate}
  \item بين أن العلاقة التراجعية عند \(p^r\) تكافئ دالة التوليد
  المستعملة في برهان حاصل أويلر.
  \item قارن موصل شكل مااس عندما \(|t|\gg|r|\) وعندما \(t\approx r\).
  \item اشرح لماذا لا يكفي حد من الشكل
  \((1+|t|+|r|)^{1/2-\eta}\) لإثبات دون التحدب في كل مناطق هبوط
  الموصل.
  \item اكتب اعتماد الثابت الضمني كاملًا في صيغة Michel--Venkatesh.
\end{enumerate}

\subsection{المستوى الثالث}
\begin{enumerate}
  \item تتبع الفرق بين العامل غير المتشعب عند المستوى \(1\) وعامل أولي
  يقسم مستوى شكل جديد، مستعينًا بمرجع خارجي مع التصريح بالتطبيع.
  \item صغ عائلة في جانب المستوى، ثم اكتب بدقة ما يعنيه حد التحدب وحد
  دون التحدب وحد ليندلوف فيها.
  \item ارسم قاموسًا محليًا--عالميًا يبدأ من \(\pi_p\) وينتهي عند
  المعادلة الوظيفية العالمية، مع تمييز المدخلات المقتبسة.
  \item اختر مثالًا مثبتًا لعملية نقل دالي خاصة، وحدد بدقة ما لا يثبته
  عن \textenglish{functoriality} العامة.
\end{enumerate}

\section*{حالة نتائج المسودة}

\begin{center}
\begin{tabular}{ll}
\textenglish{\texttt{ANT-DEF-21-01}} & \textenglish{\texttt{AUTHORED-DRAFT / NON-CITABLE}}\\
\textenglish{\texttt{ANT-PROP-21-01}} & \textenglish{\texttt{AUTHORED-DRAFT / NON-CITABLE}}\\
\textenglish{\texttt{ANT-DEF-21-02}} & \textenglish{\texttt{AUTHORED-DRAFT / NON-CITABLE}}\\
\textenglish{\texttt{ANT-THM-21-01}} & \textenglish{\texttt{AUTHORED-DRAFT / NON-CITABLE}}\\
\textenglish{\texttt{ANT-DEF-21-03}} & \textenglish{\texttt{AUTHORED-DRAFT / NON-CITABLE}}\\
\textenglish{\texttt{ANT-THM-21-02}} & \textenglish{\texttt{AUTHORED-DRAFT / NON-CITABLE}}\\
\textenglish{\texttt{ANT-DEF-21-04}} & \textenglish{\texttt{AUTHORED-DRAFT / NON-CITABLE}}\\
\textenglish{\texttt{ANT-THM-21-03}} & \textenglish{\texttt{AUTHORED-DRAFT / NON-CITABLE}}\\
\textenglish{\texttt{ANT-DEF-21-05}} & \textenglish{\texttt{AUTHORED-DRAFT / NON-CITABLE}}\\
\textenglish{\texttt{ANT-OPEN-21-01}} & \textenglish{\texttt{AUTHORED-DRAFT / NON-CITABLE}}
\end{tabular}
\end{center}

المعرفات العشرة موجودة في المتن، لكنها لا تزال غير قابلة للاستشهاد.
